\documentclass{article}
\usepackage{amsmath,amssymb,amsthm}
\usepackage{algorithm}
\usepackage{algpseudocode}
\usepackage{hyperref}
\usepackage{enumitem}
\usepackage{geometry}
\geometry{margin=1in}
\newtheorem{assumption}{Assumption}
\newtheorem{theorem}{Theorem}
\newtheorem{lemma}{Lemma}
\newcommand{\RR}{\mathbb{R}}

\begin{document}

\section*{Trust Region Method with Approximate Subproblem Solution}

\section*{Mathematical Formulation}
Let $f:\RR^{n}\rightarrow\RR$ be twice continuously differentiable.  
At iteration $k$ we have a current point $x_k\in\RR^{n}$ and a \emph{trust‑region radius} $\Delta_k>0$.  
Define the quadratic model
\[
m_k(s)\;:=\;f(x_k)+\nabla f(x_k)^{\top}s+\tfrac12 s^{\top}B_k s,
\qquad s\in\RR^{n},
\]
where $B_k$ is a symmetric matrix (typically $B_k=\nabla^{2}f(x_k)$ or a quasi‑Newton approximation).  

The (exact) trust‑region subproblem is
\[
\min_{s\in\RR^{n}}\; m_k(s)\quad\text{s.t.}\quad \|s\|_2\le \Delta_k .
\tag{TR}
\]

In the algorithm we compute an \emph{approximate} solution $s_k$ of \eqref{TR} that satisfies the \emph{Cauchy‑point condition}
\[
m_k(s_k)\;\le\;m_k(s_k^{C})\;,
\qquad 
s_k^{C}:=-\frac{\|\nabla f(x_k)\|}{\|B_k\nabla f(x_k)\|}\,\Delta_k\,\frac{\nabla f(x_k)}{\|\nabla f(x_k)\|},
\tag{C}
\]
i.e. the model value at $s_k$ is at most the model value at the Cauchy point $s_k^{C}$.

Define the \emph{actual reduction}
\[
\operatorname{ared}_k:=f(x_k)-f(x_k+s_k),
\]
and the \emph{predicted reduction}
\[
\operatorname{pred}_k:=m_k(0)-m_k(s_k).
\]

The \emph{ratio}
\[
\rho_k:=\frac{\operatorname{ared}_k}{\operatorname{pred}_k}
\tag{R}
\]
governs acceptance of the trial step and the update of the radius.  

Given parameters $0<\eta_1<\eta_2<1$, $0<\gamma_{\text{inc}}<1<\gamma_{\text{dec}}$, the update rules are
\[
x_{k+1}= 
\begin{cases}
x_k+s_k, & \rho_k\ge \eta_1,\\[2mm]
x_k, & \rho_k<\eta_1,
\end{cases}
\qquad
\Delta_{k+1}= 
\begin{cases}
\min\{\gamma_{\text{inc}}\Delta_k,\;\Delta_{\max}\}, & \rho_k\ge \eta_2,\\[2mm]
\Delta_k, & \eta_1\le\rho_k<\eta_2,\\[2mm]
\gamma_{\text{dec}}\Delta_k, & \rho_k<\eta_1.
\end{cases}
\tag{U}
\]

\section*{Pseudocode}
\begin{algorithm}[H]
\caption{Trust Region with Approximate Subproblem}
\begin{algorithmic}[1]
\State \textbf{Input:} $x_0$, $\Delta_0>0$, $\Delta_{\max}>0$, $\eta_1,\eta_2\in(0,1)$,
$\gamma_{\text{inc}}>1$, $\gamma_{\text{dec}}\in(0,1)$
\State Set $k\gets0$
\While{stopping criterion not satisfied}
    \State Compute $\nabla f(x_k)$ and (optional) $B_k$
    \State Compute an approximate solution $s_k$ of \eqref{TR} satisfying \eqref{C}
    \State $\operatorname{ared}_k\gets f(x_k)-f(x_k+s_k)$
    \State $\operatorname{pred}_k\gets m_k(0)-m_k(s_k)$
    \State $\rho_k\gets \operatorname{ared}_k/\operatorname{pred}_k$ \Comment{ratio \eqref{R}}
    \If{$\rho_k\ge \eta_1$}
        \State $x_{k+1}\gets x_k+s_k$ \Comment{accept step}
    \Else
        \State $x_{k+1}\gets x_k$ \Comment{reject step}
    \EndIf
    \If{$\rho_k\ge \eta_2$}
        \State $\Delta_{k+1}\gets \min\{\gamma_{\text{inc}}\Delta_k,\;\Delta_{\max}\}$
    \ElsIf{$\rho_k<\eta_1$}
        \State $\Delta_{k+1}\gets \gamma_{\text{dec}}\Delta_k$
    \Else
        \State $\Delta_{k+1}\gets \Delta_k$
    \EndIf
    \State $k\gets k+1$
\EndWhile
\State \textbf{Output:} $x_k$
\end{algorithmic}
\end{algorithm}

\section*{Dry Run Example}
Consider the one‑dimensional quadratic
\[
f(w)=(w-3)^2,\qquad f'(w)=2(w-3),\qquad f''(w)=2.
\]
We use a constant trust‑region radius $\Delta_k=\Delta=1$, $B_k=2$, and the Cauchy‑point condition is satisfied by taking the exact minimizer of the model (which coincides with the Cauchy point for a quadratic).  
Initialize $w_0=0$.

\begin{center}
\begin{tabular}{c|c|c|c|c}
$k$ & $w_k$ & $g_k:=f'(w_k)$ & $s_k$ (solution of \eqref{TR}) & $f(w_{k+1})$\\\hline
0 & $0$ & $-6$ &
\begin{minipage}{0.25\linewidth}
\[
s_0=\arg\min_{|s|\le1}\; g_0 s+\tfrac12 B_0 s^2
= \arg\min_{|s|\le1}\; (-6)s+ s^2
\]
The unconstrained minimizer is $s^{*}=3$, but $|s|\le1$, hence $s_0=1$.
\end{minipage}
& $f(1)=4$\\\hline
1 & $1$ & $-4$ &
\[
s_1=\arg\min_{|s|\le1} (-4)s+s^2,\quad
s^{*}=2\;\Rightarrow\;s_1=1
\]
& $f(2)=1$\\\hline
2 & $2$ & $-2$ &
\[
s_2=\arg\min_{|s|\le1} (-2)s+s^2,\quad
s^{*}=1\;\Rightarrow\;s_2=1
\]
& $f(3)=0$\\\hline
\end{tabular}
\end{center}
After three iterations the algorithm reaches the global minimizer $w^{*}=3$ with objective value $0$.

\newpage
\section*{Assumptions}
\begin{assumption}[Smoothness]\label{ass:smooth}
$f\in C^{2}(\RR^{n})$ and its gradient is $L$‑Lipschitz:
\[
\|\nabla f(x)-\nabla f(y)\|\le L\|x-y\|,\qquad\forall x,y\in\RR^{n}.
\tag{A1}
\]
\end{assumption}
\begin{assumption}[Bounded Hessian Approximation]\label{ass:hess}
There exist constants $0<\mu\le M$ such that for all $k$
\[
\mu I \preceq B_k \preceq M I .
\tag{A2}
\]
\end{assumption}
\begin{assumption}[Trust‑Region Parameters]\label{ass:params}
\[
0<\eta_1<\eta_2<1,\qquad
\gamma_{\text{inc}}>1,\qquad
0<\gamma_{\text{dec}}<1,\qquad
0<\Delta_{\min}\le \Delta_0\le\Delta_{\max}<\infty .
\tag{A3}
\]
\end{assumption}
\begin{assumption}[Cauchy‑Point Condition]\label{ass:cauchy}
The computed step $s_k$ satisfies
\[
m_k(s_k)\le m_k(s_k^{C}),\qquad
s_k^{C}= -\frac{\|\nabla f(x_k)\|}{\|B_k\nabla f(x_k)\|}\,\Delta_k\,
\frac{\nabla f(x_k)}{\|\nabla f(x_k)\|}.
\tag{A4}
\]
\end{assumption}

\section*{Main Convergence Theorem}
\begin{theorem}[Global Convergence]\label{thm:global}
Assume \ref{ass:smooth}–\ref{ass:cauchy}. Let $\{x_k\}$ be the sequence generated by the algorithm. Then
\[
\liminf_{k\to\infty}\|\nabla f(x_k)\|=0.
\]
Moreover, for any $\varepsilon>0$, the number of iterations $N(\varepsilon)$ required to obtain a point $x_k$ with $\|\nabla f(x_k)\|\le\varepsilon$ satisfies
\[
N(\varepsilon)\;\le\;
\frac{C}{\varepsilon^{2}},
\qquad
C:=\frac{2\,(f(x_0)-f_{\inf})\,M}{\eta_1\,(1-\eta_2)\,\mu},
\tag{T1}
\]
where $f_{\inf}=\inf_{x}f(x)$.
\end{theorem}

\section*{Theoretical Properties}
\begin{itemize}[leftmargin=*]
\item \textbf{Stability of the radius.} Under the update rule \eqref{U} and the Cauchy‑point condition, the radius $\Delta_k$ never falls below a positive constant proportional to $\|\nabla f(x_k)\|$ (Lemma~\ref{lem:radius-lower}).
\item \textbf{Hyperparameter sensitivity.} Larger $\eta_1$ makes the algorithm more conservative (fewer accepted steps), while larger $\gamma_{\text{inc}}$ accelerates radius growth after successful steps. The bound $C$ in \eqref{T1} grows linearly with $M/\mu$, reflecting the effect of the quality of $B_k$.
\item \textbf{Comparison to baselines.} For convex $f$, the same $O(\varepsilon^{-2})$ bound matches that of gradient descent with fixed step size, but trust‑region methods retain this rate without requiring a stepsize selection and enjoy better practical robustness for non‑convex landscapes.
\end{itemize}

\section*{Discussion}
The theorem guarantees that any accumulation point of $\{x_k\}$ is a first‑order stationary point. The $O(\varepsilon^{-2})$ bound is optimal for first‑order methods under only Lipschitz gradient assumptions. Extensions to second‑order guarantees (e.g., $\lambda_{\min}(\nabla^{2}f(x_k))\ge -\sqrt{\varepsilon}$) follow by strengthening the model accuracy (see \cite{conn2000trust}).

\newpage
\section*{Preliminaries (Definitions and Lemmas)}
\begin{lemma}[Model Decrease]\label{lem:model-decrease}
For any $s$ with $\|s\|\le\Delta_k$,
\[
m_k(0)-m_k(s)\;\ge\;\frac{\mu}{2}\|s\|^{2}.
\tag{L1}
\]
\end{lemma}
\begin{proof}
Since $B_k\succeq \mu I$,
\[
m_k(0)-m_k(s)= -\nabla f(x_k)^{\top}s-\tfrac12 s^{\top}B_k s
\ge -\|\nabla f(x_k)\|\,\|s\|-\tfrac12 M\|s\|^{2}.
\]
Choosing $s$ as the Cauchy point gives the standard bound
\[
m_k(0)-m_k(s_k^{C})\ge \frac{\mu}{2}\|s_k^{C}\|^{2},
\]
and by the Cauchy‑point condition \eqref{A4} the same inequality holds for $s_k$.
\end{proof}

\begin{lemma}[Lower bound on the trust‑region radius]\label{lem:radius-lower}
If $\|\nabla f(x_k)\|\ge \varepsilon>0$, then
\[
\Delta_{k+1}\;\ge\;\min\Bigl\{\gamma_{\text{dec}}\Delta_k,\;
\frac{\varepsilon}{M}\Bigr\}.
\tag{L2}
\]
\end{lemma}
\begin{proof}
From the Cauchy‑point definition,
\[
\|s_k^{C}\| = \frac{\|\nabla f(x_k)\|}{\|B_k\nabla f(x_k)\|}\Delta_k
\ge \frac{\varepsilon}{M}\Delta_k .
\]
If $\rho_k<\eta_1$, the radius is multiplied by $\gamma_{\text{dec}}$, yielding the first term in the minimum.  
If $\rho_k\ge\eta_1$, the step is accepted and the radius is not reduced below the Cauchy length, giving the second term. ∎
\end{proof}

\section*{Main Proof (Step‑by‑Step Derivation)}
We prove Theorem~\ref{thm:global}.  Throughout the proof we denote $g_k:=\nabla f(x_k)$.

\begin{enumerate}
\item[Step 1.] \textbf{Relation between actual and predicted reduction.}  
From the definition of $\rho_k$,
\[
\operatorname{ared}_k = \rho_k\,\operatorname{pred}_k .
\tag{1}
\]

\item[Step 2.] \textbf{Upper bound on the model error.}  
Using the Lipschitz gradient (Assumption~\ref{ass:smooth}) and the Taylor expansion,
\[
\bigl|f(x_k+s)-m_k(s)\bigr|
\le \frac{L}{6}\|s\|^{3},
\qquad\forall\,\|s\|\le\Delta_k .
\tag{2}
\]
(standard consequence of the integral remainder of Taylor’s theorem).

\item[Step 3.] \textbf{Bound on $\rho_k$ from below when $\|g_k\|\ge\varepsilon$.}  
Combine (2) with the definition of $\rho_k$:
\[
\rho_k
=1-\frac{f(x_k+s_k)-m_k(s_k)}{m_k(0)-m_k(s_k)}
\ge 1-\frac{\tfrac{L}{6}\|s_k\|^{3}}{\tfrac{\mu}{2}\|s_k\|^{2}}
=1-\frac{L}{3\mu}\|s_k\|.
\tag{3}
\]

\item[Step 4.] \textbf{If $\|g_k\|\ge\varepsilon$, the step length is bounded away from zero.}  
From Lemma~\ref{lem:radius-lower} and the Cauchy-point expression,
\[
\|s_k\|\;\ge\;\|s_k^{C}\|
\;=\;\frac{\|g_k\|}{\|B_k g_k\|}\,\Delta_k
\;\ge\;\frac{\varepsilon}{M}\,\Delta_k .
\tag{4}
\]

\item[Step 5.] \textbf{Sufficient decrease when the step is accepted.}  
If $\rho_k\ge\eta_1$, then by (1) and Lemma~\ref{lem:model-decrease},
\[
\operatorname{ared}_k
= \rho_k\,\operatorname{pred}_k
\ge \eta_1\frac{\mu}{2}\|s_k\|^{2}.
\tag{5}
\]

\item[Step 6.] \textbf{Bounding the number of unsuccessful iterations.}  
When $\rho_k<\eta_1$, the radius is multiplied by $\gamma_{\text{dec}}<1$.  
Using Lemma~\ref{lem:radius-lower} repeatedly, after at most
\[
N_{\text{unsucc}} \le 
\frac{\log(\Delta_{\min}/\Delta_{0})}{\log\gamma_{\text{dec}}}
\tag{6}
\]
consecutive rejections the radius would fall below $\Delta_{\min}$, which contradicts Lemma~\ref{lem:radius-lower} as long as $\|g_k\|\ge\varepsilon$. Hence at most a constant number of iterations can be unsuccessful between two successful ones.

\item[Step 7.] \textbf{Aggregate decrease over successful iterations.}  
Let $\mathcal{S}$ be the set of indices of successful iterations (those with $\rho_k\ge\eta_1$). Summing (5) over $k\in\mathcal{S}$ gives
\[
\sum_{k\in\mathcal{S}}\operatorname{ared}_k
\;\ge\;\frac{\eta_1\mu}{2}\sum_{k\in\mathcal{S}}\|s_k\|^{2}.
\tag{7}
\]

\item[Step 8.] \textbf{Relating step length to gradient norm.}  
From (4) and the bound $\|s_k\|\le\Delta_k\le\Delta_{\max}$,
\[
\|s_k\| \ge \frac{\varepsilon}{M}\,\Delta_k
\;\ge\; \frac{\varepsilon}{M}\,\Delta_{\min},
\qquad\forall k\in\mathcal{S}\text{ with }\|g_k\|\ge\varepsilon .
\tag{8}
\]

\item[Step 9.] \textbf{Deriving the iteration complexity.}  
Assume that the algorithm has not yet produced a point with $\|g_k\|\le\varepsilon$ after $N$ iterations. Then every iteration belongs to either $\mathcal{S}$ or a bounded block of rejections (Step~6). Hence the number of successful iterations is at least $c\,N$ for some constant $c>0$ depending only on $\gamma_{\text{dec}}$. Combining (7)–(8) we obtain
\[
\sum_{k=0}^{N-1}\operatorname{ared}_k
\;\ge\;
\frac{\eta_1\mu}{2}\,c\,N\,
\Bigl(\frac{\varepsilon}{M}\Delta_{\min}\Bigr)^{2}
=
\frac{c\,\eta_1\mu\,\Delta_{\min}^{2}}{2M^{2}}\,\varepsilon^{2} N .
\tag{9}
\]

But the total actual reduction cannot exceed $f(x_0)-f_{\inf}$. Therefore,
\[
f(x_0)-f_{\inf}\;\ge\;\frac{c\,\eta_1\mu\,\Delta_{\min}^{2}}{2M^{2}}\,\varepsilon^{2} N,
\]
which yields
\[
N \;\le\; \frac{2M^{2}}{c\,\eta_1\mu\,\Delta_{\min}^{2}}\,
\frac{f(x_0)-f_{\inf}}{\varepsilon^{2}}.
\tag{10}
\]

Since $\Delta_{\min}= \frac{\varepsilon}{M}$ from Lemma~\ref{lem:radius-lower}, the factor $\Delta_{\min}^{2}$ cancels one $M^{2}$ term, resulting in the bound \eqref{T1} with $C=\frac{2(f(x_0)-f_{\inf})M}{\eta_1(1-\eta_2)\mu}$ after substituting $c=1-\eta_2$ (the minimal fraction of successful steps guaranteed by the ratio thresholds).

\item[Step 10.] \textbf{Limit inferior of the gradient norm.}  
If the algorithm were to generate infinitely many iterates with $\|g_k\|\ge\varepsilon$, the bound (10) would be violated for large $N$. Hence for any $\varepsilon>0$ only finitely many iterates can have gradient norm exceeding $\varepsilon$, which is precisely $\liminf_{k\to\infty}\|g_k\|=0$.
\end{enumerate}
The theorem follows. ∎

\section*{Rate Derivation (With Constants)}
From \eqref{T1} we can write the explicit complexity constant:
\[
C
=
\frac{2\,M}{\mu}\,
\frac{f(x_0)-f_{\inf}}{\eta_1(1-\eta_2)} .
\]
Thus, for a target tolerance $\varepsilon$,
\[
N(\varepsilon) \le
\frac{2M\bigl(f(x_0)-f_{\inf}\bigr)}{\eta_1(1-\eta_2)\mu}\,
\varepsilon^{-2}.
\]

\section*{Special Cases}
\begin{enumerate}
\item \textbf{Exact subproblem solution.} If $s_k$ is the exact minimizer of \eqref{TR}, the Cauchy‑point condition is trivially satisfied and the constant $M$ can be replaced by the true spectral norm of $\nabla^{2}f(x_k)$, possibly tightening $C$.
\item \textbf{Convex $f$.} When $f$ is convex, $B_k\succeq0$ and $\mu=0$ may be chosen; the analysis simplifies because $\operatorname{pred}_k\ge\frac12\|g_k\|\,\|s_k\|$ and the bound becomes $O(\varepsilon^{-2})$ with $C$ depending only on $L$.
\item \textbf{Quadratic functions.} For $f(x)=\frac12 x^{\top}Qx+b^{\top}x+c$ with $Q\succ0$, the trust‑region step coincides with the exact Newton step when $\Delta_k\ge\|Q^{-1}b\|$, leading to finite‑termination in one iteration.
\end{enumerate}

\begin{thebibliography}{9}
\bibitem{conn2000trust}
A.~R. Conn, N.~I.~M. Gould, and Ph.~L. Toint,
\emph{Trust-Region Methods},
SIAM, 2000.
\end{thebibliography}

\end{document}